% ===================================================================
%  preamble.tex — macros common to both transcriptions.
%
%  This file contains NO measurement. Every measurement is in
%  calibrate-io.tex and calibrate-fr.tex, one file per facsimile: the
%  two booklets are not of the same format, are not set at the same
%  size, and were not scanned the same way. Load the calibration BEFORE
%  this file; this one uses it and never redefines it.
%
%  The method is that of the « Kompleta Gramatiko » project:
%  ONE LINE OF THE FACSIMILE = ONE LINE OF THE PDF, same pagination,
%  same folios, same word breaks. No hyphenation is decided by TeX;
%  every one is read off the scan and written into the source.
% ===================================================================

% -------------------------------------------------------------------
% 1. PACKAGES
% -------------------------------------------------------------------
\usepackage[T1]{fontenc}
\usepackage[utf8]{inputenc}
\usepackage{XCharter}
\usepackage{geometry}
\usepackage{graphicx}
\usepackage{tikz}
\usepackage{calc}
% FONT EXPANSION, to fill what the size does not fill. Even at the
% largest size that fits, a page keeps lines shorter than the others:
% the size is bounded by the LONGEST line of the page, and the others
% stretch. The compositor of 1926 settled that with letterspacing -- he
% opened out his slack lines. The modern equivalent is expansion: each
% glyph stretches or tightens by at most 2.5 %, which absorbs the
% shortfall without touching the word spaces, hence without making them
% gape. The line breaks do not move -- every one is forced by the
% transcription.
\usepackage[tracking=true,letterspace=0,
            expansion=true,stretch=25,shrink=25]{microtype}

\geometry{
  paperwidth=\VUpapierLargeur, paperheight=\VUpapierHauteur,
  textwidth=\VUtexteLargeur,   textheight=\VUtexteHauteur,
  top=\VUmargeSup, includehead=false, includefoot=false,
  hcentering, % les deux livrets ont des marges laterales egales
}

% -------------------------------------------------------------------
% 2. EXPLICIT LINE BREAKS
% -------------------------------------------------------------------
% \nl : end of line WITHOUT a hyphen.
% \cc : end of line WITH a hyphen (the hyphen is set).
% We never use \\ : it cancels the justification of the line.
\newcommand{\nl}{\penalty-10000\relax}
\newcommand{\cc}{-\penalty-10000\relax}

% \parplein ENDS a paragraph whose last line runs to the margin — the
% case of a paragraph that carries on onto the next page.
% \parfillskip is GLOBAL: laid once, it would justify the last line of
% every following paragraph. We shut it in a group.
\newcommand{\parplein}{{\parfillskip=0pt\relax\par}}
% \ccplein : the page ends on a BROKEN WORD, and the paragraph
% continues on the next leaf. It is \cc and \parplein at once, and
% neither will do alone:
%   -- \cc lays « -\penalty-10000 ». At the end of a paragraph, that
%      penalty opens one last EMPTY line, which takes one more row on
%      the page: the composed page gains a line the facsimile has not,
%      and everything after it drops by one.
%   -- \parplein alone does not set the hyphen.
% We therefore write the hyphen, then end the paragraph at the margin,
% with no penalty. Three pages of the Ido booklet end so (leaves 66,
% 103, 104); check 5 had reported them as « \VUcontinue without
% \parplein », which was the symptom and not the cause.
\newcommand{\ccplein}{-{\parfillskip=0pt\relax\par}}

% \VUcontinue OPENS a paragraph that continues the previous page: no
% indent, the line resumes flush left.
\newcommand{\VUcontinue}{\noindent}

% THE WORD SPACES ARE NON-BREAKING.
% Without that the model leaks in silence: \nl does force a break at the
% end of the line, but does not stop TeX putting one BEFORE it if the
% matter does not fit. A line 1 % too wide would therefore re-break
% itself, without even producing an « Overfull \hbox ». With the spaces
% made non-breaking, the only legal breakpoint is the one the
% transcription placed: a line too wide becomes an « Overfull \hbox »,
% which check 1 inventories. The fault passes from silence to noise.
% The spaces keep all their elasticity: the justification is intact.
\makeatletter\newcommand{\VUespaceMot}{%
  \ifdim\spaceskip=0pt\relax
    \@tempdimb=\f@size pt
    \hskip\VUespaceRatio\@tempdimb
      plus\VUespaceEtire\@tempdimb minus\VUespaceSerre\@tempdimb
  \else \hskip\spaceskip \fi}\makeatother
\begingroup
  \catcode`\ =\active
  \gdef\VUespacesInsecables{\catcode`\ =\active
    \def {\ifhmode\ifnum\lastnodetype>-1\relax\penalty10000\VUespaceMot\fi\fi}}%
\endgroup

% THE BASELINE GRID IS IMPOSED.
%
% TeX holds its grid only as long as the lines fit in it: if the natural
% height of a line exceeds \baselineskip, it abandons the regular pitch
% and separates the lines by \lineskip, that is, by their real bulk. The
% pitch then becomes a consequence of the text, and no longer a setting.
%
% That is exactly what was happening here. The facsimile is set almost
% solid -- size 10.93 pt on a pitch of 11.68 -- and the per-page size
% rises to 12.43 pt to fill the measure. Beyond ~11.7 pt the line no
% longer fits the pitch, TeX lets go of the grid, and the pages lengthen
% with nothing to say so: twenty-nine of them overflowed the foot of the
% paper, and changing the leading by three points did almost nothing --
% the measured slope was 0.5 where it should have been 40.
%
% With \lineskiplimit at -1000pt the condition is never met and the
% pitch stays the one that was measured, whatever the size. That is what
% metal does: the baselines are at a fixed distance, imposed by the body
% of the leads, and a slightly tall letter does not push its neighbour.
%
% THE SETTING IS LAID IN THE PAGE, NOT IN THE PREAMBLE. Every LaTeX box
% calls \@parboxrestore, which calls \normalbaselines, which puts
% \lineskiplimit back to zero. Laid once and for all at the head of the
% document, the setting was therefore cancelled at the opening of every
% minipage -- that is, on each of the two hundred and ten pages. We lay
% it again in the VUpage environment, after the box is opened.

% No hyphenation decided by TeX: every break is read off the facsimile.
\hyphenpenalty=10000
\exhyphenpenalty=10000
\pretolerance=-1
\tolerance=9999
\emergencystretch=0pt

% -------------------------------------------------------------------
% 3. THE PAGE IS FROZEN
% -------------------------------------------------------------------
% Each page of the facsimile is set in a minipage of fixed height: a
% retouch in one paragraph cannot make text migrate to the next page.
\pagestyle{empty}
\setlength{\parskip}{0pt}
\setlength{\topskip}{0pt}
\raggedbottom
\setlength{\parindent}{\VUalinea}

% THE HIDDEN MARK, on every page. Set in RENDERING MODE 3 (« 3 Tr »):
% the text enters the PDF's flow, it is indexed, searchable and
% extracted by pdftotext, but it is NOT PAINTED. It is not white on
% white, which would show when selecting: it is text the rendering
% engine skips. Size 0.4 pt and a box of zero height AND width, for
% good measure.
\newcommand{\VUmarqueCachee}{%
  \vbox to 0pt{%
    \hbox to 0pt{%
      \pdfliteral{3 Tr}%
      \fontsize{0.4pt}{0.4pt}\selectfont\VUmarque
      \pdfliteral{0 Tr}\hss}%
    \vss}%
  \nointerlineskip}

% #1 (optional) = the facsimile's LEAF number. It is not set: it serves
% the checks in finding the page of the scan again.
% #2 = the printed folio; empty if the page carries none.
% -------------------------------------------------------------------
% 3 bis. THE SIZE, PAGE BY PAGE
% -------------------------------------------------------------------
% ONE SIZE FOR THE WHOLE VOLUME CANNOT FILL EVERY PAGE. The global size
% is the one that, on average, fills best; on the pages where the
% original is set a little tighter it is therefore too small -- and the
% lines stretch there, with gaping word spaces. The compositor of 1926
% did not hold a constant set from one gathering to the next, and
% today's font has not exactly his. We therefore lay a size PER PAGE,
% measured by tools/body.py: the largest that makes no line overflow.
\makeatletter
\newcommand{\VUkorpoPage}[2]{%
  \expandafter\gdef\csname VUkorpo@#1\endcsname{#2}}
\newcommand{\VUkorpoDe}[1]{%
  \ifcsname VUkorpo@#1\endcsname\csname VUkorpo@#1\endcsname\else
  \VUcorps\fi}
\newcommand{\VUtextePage}[1]{%
  \edef\VU@tmp{\VUkorpoDe{#1}}%
  \fontsize{\VU@tmp}{\VUinterlignoDe{#1}}\selectfont}
% THE LEADING IS THE PAGE'S OWN TOO. The measured pitch is a median;
% the pages that carry more lines than the median do not fit the block
% and overflow at the foot. The compositor « carded » his pages -- he
% tightened or opened his leads according to the matter. tools/height.py
% measures, page by page, the pitch that makes the matter fit the
% surveyed height.
\newcommand{\VUinterlignoPage}[2]{%
  \expandafter\gdef\csname VUintl@#1\endcsname{#2}}
\newcommand{\VUinterlignoDe}[1]{%
  \ifcsname VUintl@#1\endcsname\csname VUintl@#1\endcsname\else
  \VUinterligne\fi}
\makeatother

% THE PAGE IS SET IN A BOX THAT IS MEASURED. The fixed-height minipage
% ends with a \vss, whose shrink is INFINITE: a page too full therefore
% produces no « Overfull \vbox », it overflows the foot of the paper in
% silence. Four pages of the Ido booklet and three of the French came
% off the guillotine that way with nothing to report it, and it was the
% eye that found them.
%
% We therefore set first in a box of natural height, announce its height
% in the log -- it is tools/height.py that reads it -- then lay it in
% the page. Nothing changes in the composition; what changes is that we
% now know what we are composing.
\newsavebox{\VUpagoBoxo}
\newenvironment{VUpage}[2][]%
  {\par
   % THE \noindent GOES BEFORE THE BOX WHEN IT IS LAID, NOT BEFORE THE BOX
% WHEN IT IS MEASURED. Placed here, it opened a paragraph, the lrbox
% inserted itself into it, and the page's vbox came and swelled that
% same line: a hundred and eleven pages out of a hundred and fourteen
% returned an « Overfull \hbox » of 11.3 pt, always the same one, which
% came from no text at all.
   \begin{lrbox}{\VUpagoBoxo}%
   % THE HEIGHT IS FIXED HERE, in the measured box itself. Set later, at
% the moment of laying the page, it brought back the 111 « Overfull
% \hbox » of the indent. Set here, the box comes out at the right height
% already and is laid with no wrapper.
% The price: \typeout{VUPAGO} then always returns \textheight. To
% measure the natural heights, tools/height.py temporarily removes the
% « [\textheight] » below.
   \begin{minipage}[t][\textheight][t]{\textwidth}%
     % minipage puts \parindent back to zero (\@parboxrestore): we restore
% it, both facsimiles indent every one of their paragraphs.
     \setlength{\parindent}{\VUalinea}%
     \lineskiplimit=-1000pt \lineskip=0pt
     \VUespacesInsecables
     \VUnotofalse
     \VUtextePage{#1}\VUmarqueCachee\VUfolio{#2}\ignorespaces}%
  {\VUnotoPose
   \end{minipage}%
   \end{lrbox}%
   % HEIGHT + DEPTH. A minipage aligned by the top has its reference
% point on the FIRST baseline: nearly all its matter is therefore in
% depth, and \ht alone returned 0.0 pt for all hundred and fourteen
% pages.
   \typeout{VUPAGO \the\dimexpr\ht\VUpagoBoxo+\dp\VUpagoBoxo\relax}%
   % THE BOX IS LAID AS IT STANDS, without being set to \textheight.
% The \vbox to \textheight that wrapped it added 11.3 pt to the line
% that carried it, on each of the two hundred and ten pages -- one
% « Overfull \hbox » per page, which drowned the real ones. It served
% only to stop the matter migrating to the next page, and the
% \penalty-10000 that follows stops that already. The page begins at
% \VUmargeSup in any case, which geometry fixes.
% THE \noindent IS INDISPENSABLE. Without it the page receives the
% paragraph indent -- 11.29575 pt, exactly the 5 % of the measure that
% \VUalinea is worth -- and each of the two hundred and ten pages
% returned an « Overfull \hbox » of that width, always the same, which
% came from no text and drowned the real ones.
% AND THE HEIGHT STAYS FIXED. Laid at its natural height, a page too
% full spilled onto one page more, and the volume gained one: the rule
% « one page of the facsimile = one page of the PDF » was broken. The
% box is therefore set to \textheight, which keeps it in its place; the
% few pages that do not fit are reported by tools/height.py, and it is
% their transcription that must be looked at again.
   \noindent\usebox{\VUpagoBoxo}%
   \par\penalty-10000\relax}

% Leading proper to one page — the compositor's « carding », which
% spreads the white of a page too short into uniform leads. To be
% called just after \begin{VUpage}, BEFORE any matter. With no call,
% the page keeps \VUinterligne.
\newcommand{\VUinterlignePage}[1]{\fontsize{\VUcorps}{#1}\selectfont}

% ORDINARY white between two paragraphs. It is NOT put in \parskip: at
% the top of a minipage the glue of \parskip survives and would shift
% the whole page by that much. We therefore write it, like the line
% breaks, where it belongs.
\newcommand{\VUblancAlinea}{\vskip\VUblancAlineaValeur\relax}
% Wider white at an articulation. #1 = the EXCESS measured over the
% page's regular pitch, LESS the ordinary white when it replaces it.
\newcommand{\VUblanc}[1]{\vskip #1\relax}

% -------------------------------------------------------------------
% 4. THE BODY OF THE TEXT
% -------------------------------------------------------------------
\newcommand{\VUtexte}{\fontsize{\VUcorps}{\VUinterligne}\selectfont}
% Bold: its own family, size PROPORTIONAL to the current size — it
% therefore follows the text, the notes, or any other size in which it
% appears, automatically. The bold of both booklets is a narrow semibold
% not proportioned to its roman; XCharter's bold is wide. Hence a
% distinct family and a ratio, not an absolute size.
\makeatletter
\newcommand{\VUgras}[1]{{%
  \@tempdima=\f@size pt
  \@tempdima=\VUratioGras\@tempdima
  \fontfamily{\VUfonteGras}%
  \fontsize{\@tempdima}{\f@baselineskip}\selectfont\bfseries #1}}
\makeatother
\newcommand{\VUtexteNote}{%
  \fontsize{\VUcorpsNote}{\VUinterligneNote}\selectfont\lsstyle}

% -------------------------------------------------------------------
% 5. DISPLAY LINES
% -------------------------------------------------------------------
% On these pages too, ONE line of the facsimile = ONE macro call: the
% transcription stays line for line, and the checks can verify it.
% #1 = size; #2 = letterspacing (1/1000 of an em); #3 = text.
% The first two are MEASURED: the size by the height of the capitals,
% the letterspacing by the width of the line (tools/apparat.py).
\newcommand{\VUcentre}[3]{{\centering\fontsize{#1}{#1}\selectfont
   \textls[#2]{#3}\par}}
\newcommand{\VUsaut}[1]{\vspace{#1}}

% Section title: bold letterspaced capitals, centred.
\newcommand{\VUtitre}[3]{%
  {\centering\fontfamily{\VUfonteGras}\fontsize{#1}{#1}\selectfont\bfseries
   \textls[#2]{#3}\par}}
% HEADING IN SMALL CAPITALS — « DESKRIPTO GENERALA. », « LA BONA
% LERNANTI. », « I. DESCRIPTION GÉNÉRALE. ». Neither bold nor full
% capitals: it is the roman of the running text, of which only the
% first letter of each word keeps its full capital. Both booklets use
% it the same way — it is the rank of title IMMEDIATELY BELOW that of
% the table number.
\newcommand{\VUintitulo}[2]{%
  {\centering\fontsize{#1}{#1}\selectfont\textls[#2]{\scshape}\par}}
\newcommand{\VUpk}[3]{%
  {\centering\fontsize{#1}{#1}\selectfont\textls[#2]{\scshape #3}\par}}

% Centred rule: #1 = width.
\newcommand{\VUfilet}[1]{{\centering\rule{#1}{0.4pt}\par}}
% Note rule: short, flush left.
\newcommand{\VUfiletnote}{{\noindent\rule{\VUfiletnoteLargeur}{0.4pt}\par}}

% Folio centred, framed by two en dashes, as « — 8 — ». Placed
% ABSOLUTELY, at zero height: it pushes nothing. The body of the text
% therefore always begins exactly at \VUmargeSup, whether the page
% carries a folio or not.
\newcommand{\VUfolio}[1]{%
  \ifx\relax#1\relax\else
    \vbox to 0pt{\vskip\dimexpr\VUfolioY-\VUmargeSup\relax
      \hbox to \textwidth{\hss--\,#1\,--\hss}\vss}%
    \nointerlineskip
  \fi}

% -------------------------------------------------------------------
% 7 bis. FOOTNOTES
% -------------------------------------------------------------------
% THEY ARE SET FROM THE FOOT, AND NOT FROM THE TOP.
%
% First version: each call gave the ordinate of the RULE, and the
% transcription brief suggested 143.2 mm « if the note is at the foot of
% a full page ». All thirteen transcriptions copied that value out --
% that is my fault, not theirs: a starting ordinate cannot suit notes of
% different lengths. A note of two lines then floated in the middle of
% the white, and a note of twelve lines -- leaf 19 of the Ido booklet --
% came off the foot of the paper, its last three lines cut away at the
% guillotine.
%
% The compositor does not do that: he sets his note ON THE FOOT of the
% block, and it is the top of the rule that rises when the note grows
% longer. We therefore measure the note in a box, and lay it so that its
% last row falls at the foot of the block of text. The ordinate ceases
% to be a datum of the transcription: it is computed, and it is right
% for every length.
%
% #1 is KEPT but ignored: removing it would mean retouching the thirteen
% transcriptions for no gain, and the argument preserves a trace of the
% ordinate the facsimile shows, should we ever want to verify it.
% AND THE NOTE IS NOT LAID AT ZERO HEIGHT: IT IS DEFERRED.
%
% Second version: the note was laid ABSOLUTELY, in a \vbox to 0pt, set
% on the foot of the block. It fell right -- and it pushed nothing. On a
% page whose text runs down to the foot, the two therefore printed one
% ON the other: seven pages of the Ido booklet and one of the French,
% measured with mutool by comparing the size of the lines (the note is
% 0.79 times the text, the superscript of a cross-reference 0.70 -- it
% is the sizes that separate them, not the position).
%
% The compositor cannot do that: his note occupies the foot of the page,
% and the text stops above it. We therefore defer the note to the
% CLOSING of the page, where it enters the flow after the last
% paragraph. The overlap becomes impossible by construction, and not by
% a calculation that would have to be redone at every transcription.
%
% THE SPRING IS IN « fill » AND NOT IN « fil »: the minipage of \VUpage
% already carries a \vss at the foot, and two springs of the same order
% would share the white out -- the note would settle in the MIDDLE of
% what is left. One order above takes everything, and the note goes down
% to the foot.
%
% #1 is KEPT but ignored: removing it would mean retouching the thirteen
% transcriptions for no gain, and the argument preserves a trace of the
% ordinate the facsimile shows, should we ever want to verify it.
\newsavebox{\VUnotoBoxo}
\newsavebox{\VUnotoUno}
\newif\ifVUnoto
% Foot of the block of text, counted from the top edge of the paper.
\newcommand{\VUnotoBas}{\dimexpr\VUfolioY+\VUtexteHauteur\relax}
% THE FOOT OF THE NOTE, WHEN THE BLOCK ITSELF FALLS BELOW THE PAPER.
% The French booklet's block is calibrated at 7.96 + 173.40 = 181.36 mm
% on a paper of 180: its foot is ALREADY 1.36 mm off the sheet. That is
% the error check 6 has always reported -- the frame of the French scan
% falls inside the paper, and the absolute scale drawn from it is wrong
% (README, § 4). As long as the text stops higher, it does not show; a
% note laid at the foot of the block does come off the guillotine. We
% therefore hold it 2 mm above the edge of the paper.
% The value is ZERO for the Ido booklet, whose block ends at 161.33 mm
% on a paper of 180: it is not a setting, it is a calculation, and it
% bites only where the calibration is at fault.
\newlength{\VUnotoPied}
% THE FINAL \relax CLOSES THE ASSIGNMENT, AND ONE IS NEEDED HERE TOO:
% without it, TeX looked for « plus » beyond the \dimexpr and expanded
% the \ifdim of the next line to find it -- which was therefore
% evaluated BEFORE the register had its value, on the zero of the
% \newlength. The condition was always false, the foot stayed negative,
% and the PDF came out identical to the byte: nothing distinguishes a
% condition that does not bite from a condition that is absent.
\VUnotoPied=\dimexpr\VUmargeSup+\VUtexteHauteur-\VUpapierHauteur+2mm\relax\relax
% THE \relax AFTER « 0pt » IS NOT DECORATIVE. \VUnotoPied is a GLUE
% register: after the dimension, TeX looks for « plus » and « minus »,
% and to find them it EXPANDS what follows -- hence the \fi, which
% closes the condition before the assignment is made. Without that
% \relax, the Ido booklet kept a foot of -47.43 pt and its note climbed
% back into the text: the condition was right, the assignment did not
% reach it.
\ifdim\VUnotoPied<0pt\relax \VUnotoPied=0pt\relax \fi
% AND « \VUnotoBas » STARTED FROM THE FOLIO, WHICH IS ABOVE THE BLOCK.
% \VUfolioY gives the baseline of the leaf number, \VUmargeSup the top
% of the block of text: the two differ by 4.11 mm in Ido and 3.75 mm in
% French. Setting the note on \VUfolioY+\VUtexteHauteur therefore laid
% it that much TOO HIGH, that is, in the text. The foot of the block is
% \VUmargeSup+\VUtexteHauteur, and that is the bottom of the minipage:
% the note settles there of itself, with no correction. \VUnotoBas is
% now kept for the record only.
\newcommand{\VUnotes}[2]{%
  % \vtop AND NOT \vbox: the reference point of a \vbox is its LAST
% baseline, so that added after the text it rises by its whole height
% and settles ON TOP of it. A \vtop's is its first line: the note goes
% down, as one expects.
  \sbox{\VUnotoUno}{\vtop{\hsize=\textwidth
      \VUfiletnote\vskip\VUblancNote\VUtexteNote
      \setlength{\parindent}{\VUalineaNote}#2\par}}%
  % TWO NOTES ON ONE PAGE does happen -- leaf 41 of the Ido booklet: the
% second stacks beneath the first, they do not replace each other.
  \ifVUnoto
    \sbox{\VUnotoBoxo}{\vtop{\hbox{\usebox{\VUnotoBoxo}}\hbox{\usebox{\VUnotoUno}}}}%
  \else
    \sbox{\VUnotoBoxo}{\vtop{\hbox{\usebox{\VUnotoUno}}}}%
    \VUnototrue
  \fi
  \ignorespaces}
% What \VUpage lays at the foot of the page, if there is a note.
\newcommand{\VUnotoPose}{%
  \ifVUnoto
    \par\nobreak\vskip0pt plus1fill\relax
    % AND THE NOTE COUNTS FOR ITS WHOLE HEIGHT, NOT FOR ITS FIRST LINE.
% \VUnotoBoxo is a \vtop: its height is that of its FIRST line, all
% the rest is in DEPTH. But the height of a \vbox stops at the last
% baseline -- the note's depth therefore did not enter the
% calculation, and the spring pushed it until its first line touched
% the foot of the block. The other six lines hung below, off the leaf,
% AND THAT WHATEVER THE QUANTITY OF TEXT: leaf 17, two thirds empty,
% came off the paper exactly as leaf 9 did, full to the brim. A \kern
% equal to that depth closes the vertical list: the box has no depth
% left, its height includes the whole note, and the spring lays the
% FOOT of the note at the foot of the block -- where the compositor
% lays it.
%
% THE \par IS INDISPENSABLE BEFORE THE \kern. The \noindent opens a
% paragraph; without closing it, the \kern was laid in HORIZONTAL
% mode, where it is worth only a space of zero width. Nothing showed
% for it: the PDF came out identical to the byte.
    \noindent\hbox to \textwidth{\usebox{\VUnotoBoxo}\hss}\par
    \nobreak\kern\dimexpr\dp\VUnotoBoxo+\VUnotoPied\relax
  \fi}

% -------------------------------------------------------------------
% 7 ter. WHAT IS NOT SET IN TYPE
% -------------------------------------------------------------------
% Two things in these two volumes are not composition: the wood
% engraving of wall plate no. 12, on leaf 111 of the Ido booklet, and
% the Art Nouveau frame of the French back cover. To render them in
% type would be to invent them; to omit them — which was done at first,
% a \VUsaut and a comment — is to publish a page that lies about what it
% showed. We therefore take them off the scan (tools/images.py) and lay
% them at their surveyed size.
%
% #1 = file; #2 = width, as a fraction of the measure.
\newcommand{\VUimajo}[2]{%
  {\centering\noindent
   \includegraphics[width=#2\textwidth]{ornaments/#1}\par}}

% -------------------------------------------------------------------
% 7 quater. DISPLAY LINES AND FRAME
% -------------------------------------------------------------------
% \VUlineo : A LINE THAT DOES NOT JUSTIFY.
%
% The announcement pages and the back cover are made of short lines --
% « 1^o En feuilles separees ; » -- which the transcription noted as
% running text ended by \nl. But \nl forces a break WITHIN A JUSTIFIED
% PARAGRAPH: TeX then stretches the line from one margin to the other,
% and « 1^o En feuilles separees ; » barred the whole width, sowing
% two-em spaces between its words. The facsimile lays those lines at
% their natural set.
%
% Those lines are therefore not running text: they are display lines,
% and they are set as such -- flush left, at their own width, with no
% stretching.
% #1 = size; #2 = letterspacing; #3 = left indent; #4 = text.
\newcommand{\VUlineo}[4]{%
  {\fontsize{#1}{#1}\selectfont
   \noindent\hspace*{#3}\textls[#2]{#4}\hfill\par}}
% Bold variant, for the display slab serifs.
\newcommand{\VUlineoG}[4]{%
  {\fontfamily{\VUfonteGras}\fontsize{#1}{#1}\selectfont\bfseries
   \noindent\hspace*{#3}\textls[#2]{#4}\hfill\par}}

% \VUkadro : THE FRAME OF THE BACK COVER, drawn and not photographed.
% The facsimile carries an Art Nouveau frame: a double rule whose
% corners curve away into volutes. We do not photo\-graph it — that
% would be to give up setting the page — we draw it. The drawing below
% is not a copy of the original block: it is a frame of the same make,
% a double rule with softened corners and corner fleurons, laid at the
% surveyed dimensions. The difference is stated here rather than
% concealed.
% #1 = height of the frame.
\newcommand{\VUkadro}[2]{%
  \begingroup
  \setlength{\unitlength}{1pt}%
  \noindent\begin{tikzpicture}[baseline=(courant.north)]
    % inner sep=0pt: without it, this empty node stays framed by the
% default inner margin (0.333 em) and widens the box by 3.88 pt beyond
% \textwidth -- a constant « Overfull \hbox », identical on every page,
% whatever its content.
    \node[inner sep=0pt] (courant) at (0,0) {};
    % The outer rule is drawn 0.55 pt INSIDE each edge: a stroke of 1.1 pt
% centred on x=0 overflows by 0.55 pt to the left AND to the right of
% \textwidth, and that last overflow returned an « Overfull \hbox » of
% 1.1 pt on the whole page.
    \draw[line width=1.1pt, rounded corners=7pt]
      (0.55pt,0) rectangle (\textwidth-0.55pt, -#1);
    \draw[line width=0.4pt, rounded corners=5pt]
      (4pt,-4pt) rectangle (\textwidth-4pt, -#1+4pt);
    % The vertical offset is computed SEPARATELY, by \pgfmathsetmacro:
% passed directly in the shift={...} key, pgfmath did not evaluate it
% -- all four fleurons settled near the top, two inside, two outside,
% and the foot of the frame stayed bare.
    \foreach \x/\y/\rx/\ry in {0/0/1/-1, 1/0/-1/-1, 0/1/1/1, 1/1/-1/1} {
      \pgfmathsetmacro{\VUkadroY}{-\y*0 - (1-\y)*#1}
      \begin{scope}[shift={(\x*\textwidth, \VUkadroY pt)}]
        \draw[line width=0.9pt]
          (\rx*13pt, \ry*3pt) .. controls (\rx*7pt, \ry*10pt)
          and (\rx*3pt, \ry*7pt) .. (\rx*3pt, \ry*13pt);
      \end{scope}}
    \node[inner sep=0pt, anchor=north west] at (11pt, -13pt)
      {\begin{minipage}{\dimexpr\textwidth-22pt\relax}#2\end{minipage}};
  \end{tikzpicture}\par
  \endgroup}

% Typographic ornament (the fleuron of the table pages). The
% facsimile's drawing is not of the font: it is laid as an image once
% surveyed, and is replaced by a tilde and a space until it is.
\newcommand{\VUornamento}[1]{{\centering\fontsize{#1}{#1}\selectfont
  \textls[300]{\textasciitilde\textasciitilde\textasciitilde}\par}}

% -------------------------------------------------------------------
% 8. THE TABLE OF SIZES, PAGE BY PAGE
% -------------------------------------------------------------------
% LAST, AND THAT IS THE WHOLE POINT. This load was at first placed at
% the start of the file, before even \geometry. \VUkorpoPage was not yet
% defined there: LaTeX therefore took each call for text, and set
% « 911.60pt 1011.15pt 1111.30pt... » -- the contents of the settings
% file printed as the pages went by. The symptom was ten inexplicable
% « Overfull \hbox »; the cause was an order.
% The file may not exist: the volume is then set at the global size,
% which is the wanted behaviour before the first measurement.
\InputIfFileExists{body-\VUlangue.tex}{}{}
% AND THE TABLE IS CORRECTED AFTERWARDS, FOR THE PAGES WITH A NOTE.
% The per-page pitch above was set so that the matter FILLS the block --
% and the note did not count as matter: it is deferred to the closing of
% the page, and the measuring tool saw nothing of it. On a page with a
% note, the pitch obtained therefore fills the block with the TEXT
% ALONE, and nothing is left for the note. Leaf 9 thus received 15.31 pt
% where its neighbours hold 11.19 and 12.42, and where the facsimile
% itself measures 11.0 (inventory-io.json: a pitch of 72 px on a block
% of 2800 px for 150.65 mm).
% The file below gives those pages back a pitch that leaves room for the
% note. It is WRITTEN BY HAND, page by page: the tool that produced
% body-*.tex is no longer in the repository, and remaking it for one
% page would have cost more than measuring it.
\InputIfFileExists{body-\VUlangue-noto.tex}{}{}
